\textbf{The vehicle is a Value-at-Risk study}, chosen because it is a domain with well-established
baselines, a strictly consistent scoring function, and standard backtests --- an unusually
favourable setting for detecting a spurious result. The question was narrow: does an L2 penalty
anchoring a quantile-loss model to a classical VaR prior improve one-day-ahead forecasts? The
answer is no. That answer is not the contribution, and readers interested in VaR modelling
specifically will find \Cref{sec:limitations} disappointing by design.

A survey of \textbf{five} machine-learning VaR papers, each read in full and coded against a schema
fixed before any paper was read (\Cref{app:survey}), locates the gap more precisely. Practice in this
literature is \textbf{heterogeneous, and no paper in the sample is careless} --- two use a held-out
selection block, one selects on a strictly consistent loss, one corrects for multiplicity and
then explicitly downgrades its own p-values, and all five gate on coverage.

Across all five, however: \textbf{test-set reuse is never counted, no power analysis appears, and none
is pre-registered.} The controls this field has institutionalised are not the controls that
defeated this study. A convenience sample of five cannot establish a base rate and we do not
claim it does; by the standard \Cref{sec:power} applies to our own nulls, it is an illustration.

\textbf{Why a VaR study, and why that is not incidental.} The mechanism documented here --- an L2
penalty contracting inter-seed dispersion toward whatever target it is given --- is Stein/ridge
shrinkage, and it is not new. Neither is the control. A reader is therefore entitled to ask why
a known mechanism and a known control are worth a paper. The answer is the domain. Quantile
forecasting for VaR is unusually well instrumented: established parametric and historical
baselines, a strictly consistent scoring function, and a set of backtests that practitioners and
regulators actually run. If a mechanically guaranteed result can be produced \emph{here} --- and
survive pre-registration, a multiplicity correction, a ten-seed protocol, Diebold--Mariano tests,
a Model Confidence Set screen, a coverage gate, replication across four runs and a dose--response
relationship --- then the failure is not a field being careless with a weak checklist.
\textbf{The instrumentation is what is being tested, and a well-instrumented domain is the strongest
available place to test it.} A synthetic demonstration alone would establish that the mechanism
exists; it would not establish that the mechanism survives a real protocol built to catch
exactly this. That second claim is the contribution, and it is why \Cref{sec:multiplicity} to \Cref{sec:tautology} are in the paper
at all.

\textbf{What the market study is, and is not.} \Cref{sec:disclosure} discloses 1,959 test-set evaluations across 16
asset--level cells with \textbf{zero cells scored only once}. By this project's own protocol that means
the study has several validation blocks and no test set, and every market number in \Cref{sec:findings} is
validation-grade. We therefore do not offer \Cref{sec:multiplicity}--\Cref{sec:tautology} as evidence about VaR modelling, nor as
evidence that the mechanism operates. \textbf{They are the exhibit} --- the artefact that a careful
pre-registered protocol produced, and would have reported. The evidence that the mechanism
guarantees the result is \Cref{sec:synthetic}, where ground truth is known and the demonstration needs no market
data at all. Keeping the two apart is the single distinction this paper asks a reader to hold.

\textbf{Contributions.} Two.

\begin{enumerate}
\def\labelenumi{\arabic{enumi}.}
\tightlist
\item
  \textbf{A mechanically guaranteed result, shown to be reportable.} A demonstration --- analytical
  and synthetic, with ground truth known --- that \textbf{shrinkage-induced stability is not evidence
  of a good prior}, together with a scale-matched permutation control that separates the two
  for minutes of compute. \textbf{Neither the mechanism nor the control is new}: the mechanism is
  Stein/ridge shrinkage, and the control is a restricted randomization test with precedents in
  statistics, machine learning and econometrics (\Cref{sec:synthetic}). What is documented is that the artefact
  is \emph{reportable as an empirical finding} and survives a pre-registered protocol with a
  dose--response and \(p = 5.2 \times 10^{-4}\) --- and that we located no prior application of the
  control to a stability claim, which we state as a search null rather than as a gap.
\item
  \textbf{Four failure modes observed in a single pre-registered study}, each with the control that
  detects it and that control's cost. We claim instances, not coverage: one study cannot
  establish a taxonomy and we do not present one.
\end{enumerate}

Two further outputs are recorded rather than claimed as contributions. The study is a
pre-registered VaR null with dated amendments, several of which weakened the authors' own
position --- pre-registration remains rare in this domain, with \textbf{$\leq$ 1\% of economics journals}
having adopted the registered-report format as of July 2023 \citep{lin2024registered}. And \Cref{app:ledger}
describes an append-only \textbf{test-touch ledger} that makes evaluation reuse countable rather than
reconstructible after the fact.
